% GENERATED FILE. Edit canonical lessons, then run npm run generate:books.
% canonical-source-hash: fnv1a64:a6b8343ab0854664
% canonical-lessons: MR-C525-kalsi, MR-C525-math, MR-C525-otne, MR-C525-sandne, MR-C525-visalne

\chapter{Pouring and Rinsing}
\label{ch:mr-pouring-and-rinsing}
\hlchaptermodality{525}

\begin{chapteropening}
\textbf{By the end of this chapter:} \emph{I can talk about water pots and pouring, spilling and rinsing.}

Complete the last lesson of chapter 525: I can talk about water pots and pouring, spilling and rinsing.
\end{chapteropening}

\section[\texorpdfstring{\mr{कळशी} (kaḷśī)}{कळशी (kaḷśī)}]{\mr{कळशी} (kaḷśī) — a brass water pot}

\label{lesson:MR-C525-kalsi}



\begin{quote}
Before the new one: say the Marathi for to shake off, then the Marathi for to scrub.
\end{quote}

\subsection*{You'll want to know: \mr{कळशी}}
\textbf{\mr{कळशी}} — \emph{kaḷśī} — \textquotedblleft{}a brass water pot\textquotedblright{}.

\begin{quote}
\mr{आजीने} \mr{नळावरून} \mr{कळशी} \mr{भरून} \mr{आणली}.
\end{quote}

\begin{quote}
\emph{ājīne naḷāvarūn kaḷśī bharūn āṇlī.}
\end{quote}

\begin{quote}
Grandmother filled a water pot at the tap and brought it.
\end{quote}

Feminine.

\begin{cousinweb}
\textbf{\mr{कळशी}} (\emph{kaḷśī}) is a tadbhava, inherited from Sanskrit \emph{kalaśa}, \textquotedblleft{}waterpot\textquotedblright{} (CDIAL 2920).
\end{cousinweb}

\begin{grammarlens}[title={What you've built}]
One more word for describing what you lived and felt.
\end{grammarlens}

\subsection*{Your turn}
\begin{itemize}
\raggedright
  \item \emph{Say it:} \emph{kaḷśī}
  \item \emph{Say it:} \emph{kaḷśī}, once more
  \item \emph{Say it:} \emph{ājīne naḷāvarūn kaḷśī bharūn āṇlī.} — Grandmother filled a water pot at the tap and brought it.
  \item \emph{Recall:} say the Marathi for to shake off, then the Marathi for to scrub, then say \emph{kaḷśī} again
\end{itemize}

\begin{tcolorbox}[breakable,colback=teal!4,colframe=teal!35!black,title={Before you move on}]
What does \mr{कळशी} mean? (\textquotedblleft{}A brass water pot\textquotedblright{}.) Say it once more.
\end{tcolorbox}
\section[\texorpdfstring{\mr{माठ} (māṭh)}{माठ (māṭh)}]{\mr{माठ} (māṭh) — a clay water pot}

\label{lesson:MR-C525-math}



\begin{quote}
Before the new one: say the Marathi for to scrub, then the Marathi for a brass water pot.
\end{quote}

\subsection*{You'll want to know: \mr{माठ}}
\textbf{\mr{माठ}} — \emph{māṭh} — \textquotedblleft{}a clay water pot\textquotedblright{}.

\begin{quote}
\mr{उन्हाळ्यात} \mr{माठातलं} \mr{पाणी} \mr{थंड} \mr{राहतं}.
\end{quote}

\begin{quote}
\emph{unhāḷyāt māṭhātlaṁ pāṇī thaṇḍ rāhtaṁ.}
\end{quote}

\begin{quote}
In summer the water in the clay pot stays cool.
\end{quote}

Masculine.

\begin{cousinweb}
\textbf{\mr{माठ}} (\emph{māṭh}) is the word for the big clay pot that keeps water cool; its origin is not settled.
\end{cousinweb}

\begin{grammarlens}[title={What you've built}]
One more word for describing what you lived and felt.
\end{grammarlens}

\subsection*{Your turn}
\begin{itemize}
\raggedright
  \item \emph{Say it:} \emph{māṭh}
  \item \emph{Say it:} \emph{māṭh}, once more
  \item \emph{Say it:} \emph{unhāḷyāt māṭhātlaṁ pāṇī thaṇḍ rāhtaṁ.} — In summer the water in the clay pot stays cool.
  \item \emph{Recall:} say the Marathi for to scrub, then the Marathi for a brass water pot, then say \emph{māṭh} again
\end{itemize}

\begin{tcolorbox}[breakable,colback=teal!4,colframe=teal!35!black,title={Before you move on}]
What does \mr{माठ} mean? (\textquotedblleft{}A clay water pot\textquotedblright{}.) Say it once more.
\end{tcolorbox}
\section[\texorpdfstring{\mr{ओतणे} (otṇe)}{ओतणे (otṇe)}]{\mr{ओतणे} (otṇe) — to pour}

\label{lesson:MR-C525-otne}



\begin{quote}
Before the new one: say the Marathi for a brass water pot, then the Marathi for a clay water pot.
\end{quote}

\subsection*{You'll want to know: \mr{ओतणे}}
\textbf{\mr{ओतणे}} — \emph{otṇe} — \textquotedblleft{}to pour\textquotedblright{}.

\begin{quote}
\mr{मी} \mr{माठातलं} \mr{पाणी} \mr{कळशीत} \mr{ओतलं}.
\end{quote}

\begin{quote}
\emph{mī māṭhātlaṁ pāṇī kaḷśīt otlaṁ.}
\end{quote}

\begin{quote}
I poured the water from the clay pot into the water pot.
\end{quote}

Past, with -\mr{ने} (he …): \mr{त्याने} \mr{पाणी} \mr{ओतलं}, \emph{tyāne pāṇī otlaṁ}.

Where it goes takes -\mr{त}: \mr{कळशीत} \mr{ओतणे}.

\begin{cousinweb}
\textbf{\mr{ओतणे}} (\emph{otṇe}) is inherited: CDIAL lists Marathi \emph{otṇẽ}, \textquotedblleft{}to pour\textquotedblright{}, under \emph{*avōtta}, \textquotedblleft{}flowed down\textquotedblright{} (CDIAL 902).
\end{cousinweb}

\begin{grammarlens}[title={What you've built}]
One more word for describing what you lived and felt.
\end{grammarlens}

\subsection*{Your turn}
\begin{itemize}
\raggedright
  \item \emph{Say it:} \emph{otṇe}
  \item \emph{Say it:} \emph{otṇe}, once more
  \item \emph{Say it:} \emph{mī māṭhātlaṁ pāṇī kaḷśīt otlaṁ.} — I poured the water from the clay pot into the water pot.
  \item \emph{Recall:} say the Marathi for a brass water pot, then the Marathi for a clay water pot, then say \emph{otṇe} again
\end{itemize}

\begin{tcolorbox}[breakable,colback=teal!4,colframe=teal!35!black,title={Before you move on}]
What does \mr{ओतणे} mean? (\textquotedblleft{}To pour\textquotedblright{}.) Say it once more.
\end{tcolorbox}
\section[\texorpdfstring{\mr{सांडणे} (sāṇḍṇe)}{सांडणे (sāṇḍṇe)}]{\mr{सांडणे} (sāṇḍṇe) — to spill}

\label{lesson:MR-C525-sandne}



\begin{quote}
Before the new one: say the Marathi for a clay water pot, then the Marathi for to pour.
\end{quote}

\subsection*{You'll want to know: \mr{सांडणे}}
\textbf{\mr{सांडणे}} — \emph{sāṇḍṇe} — \textquotedblleft{}to spill\textquotedblright{}.

\begin{quote}
\mr{घाईत} \mr{माझ्याकडून} \mr{चहा} \mr{सांडला}.
\end{quote}

\begin{quote}
\emph{ghāīt mājhyākaḍūn cahā sāṇḍlā.}
\end{quote}

\begin{quote}
In my hurry I spilled the tea.
\end{quote}

Past, with -\mr{ने} (he …): \mr{त्याने} \mr{दूध} \mr{सांडलं}, \emph{tyāne dūdh sāṇḍlaṁ}.

For an accident Marathi says \mr{माझ्याकडून} \mr{चहा} \mr{सांडला}, the tea spilled from me.

\begin{cousinweb}
\textbf{\mr{सांडणे}} (\emph{sāṇḍṇe}) is inherited: CDIAL lists Marathi \emph{sā̃ḍṇẽ}, \textquotedblleft{}to spill, give up\textquotedblright{}, under Sanskrit \emph{chardati}, \textquotedblleft{}ejects, vomits\textquotedblright{} (CDIAL 4998), whose basic sense is letting something out.
\end{cousinweb}

\begin{grammarlens}[title={What you've built}]
One more word for describing what you lived and felt.
\end{grammarlens}

\subsection*{Your turn}
\begin{itemize}
\raggedright
  \item \emph{Say it:} \emph{sāṇḍṇe}
  \item \emph{Say it:} \emph{sāṇḍṇe}, once more
  \item \emph{Say it:} \emph{ghāīt mājhyākaḍūn cahā sāṇḍlā.} — In my hurry I spilled the tea.
  \item \emph{Recall:} say the Marathi for a clay water pot, then the Marathi for to pour, then say \emph{sāṇḍṇe} again
\end{itemize}

\begin{tcolorbox}[breakable,colback=teal!4,colframe=teal!35!black,title={Before you move on}]
What does \mr{सांडणे} mean? (\textquotedblleft{}To spill\textquotedblright{}.) Say it once more.
\end{tcolorbox}
\section[\texorpdfstring{\mr{विसळणे} (visaḷṇe)}{विसळणे (visaḷṇe)}]{\mr{विसळणे} (visaḷṇe) — to rinse}

\label{lesson:MR-C525-visalne}



\begin{quote}
Before the new one: say the Marathi for to pour, then the Marathi for to spill.
\end{quote}

\subsection*{You'll want to know: \mr{विसळणे}}
\textbf{\mr{विसळणे}} — \emph{visaḷṇe} — \textquotedblleft{}to rinse\textquotedblright{}.

\begin{quote}
\mr{चहा} \mr{पिऊन} \mr{झाल्यावर} \mr{त्याने} \mr{कप} \mr{विसळला}.
\end{quote}

\begin{quote}
\emph{cahā piūn jhālyāvar tyāne kap visaḷlā.}
\end{quote}

\begin{quote}
After drinking his tea he rinsed the cup.
\end{quote}

Past, with -\mr{ने} (he …): \mr{त्याने} \mr{कप} \mr{विसळला}, \emph{tyāne kap visaḷlā}.

\begin{cousinweb}
\textbf{\mr{विसळणे}} (\emph{visaḷṇe}) is inherited: CDIAL lists Marathi \emph{visaḷṇẽ}, \textquotedblleft{}to rinse\textquotedblright{}, under \emph{*vikṣalati}, \textquotedblleft{}rinses\textquotedblright{} (CDIAL 11645).
\end{cousinweb}

\begin{grammarlens}[title={What you've built}]
One more word for describing what you lived and felt.
\end{grammarlens}

\subsection*{Your turn}
\begin{itemize}
\raggedright
  \item \emph{Say it:} \emph{visaḷṇe}
  \item \emph{Say it:} \emph{visaḷṇe}, once more
  \item \emph{Say it:} \emph{cahā piūn jhālyāvar tyāne kap visaḷlā.} — After drinking his tea he rinsed the cup.
  \item \emph{Recall:} say the Marathi for to pour, then the Marathi for to spill, then say \emph{visaḷṇe} again
\end{itemize}

\begin{tcolorbox}[breakable,colback=teal!4,colframe=teal!35!black,title={Before you move on}]
What does \mr{विसळणे} mean? (\textquotedblleft{}To rinse\textquotedblright{}.) Say it once more.
\end{tcolorbox}
